\section{A uniform streaming law through an additive collision}\label{sec:uniform}
We apply the confluent law to a family in which exact dimension changes
at a fixed horizon and at a fixed positivity margin. In this family all
repeated quantization errors can be controlled in the physical prediction
metric, uniformly at the collision.

\subsection{One fixed five-trial protocol}
Let $0\le\theta\le1/2$, $A_\theta=\{0,1,2+\theta\}$, and let $\mu$
be any fixed full-support probability on $[0,1]$. Define
\begin{equation}\label{eq:near-collision-cells}
 \begin{aligned}
 k_1^\theta&=\tfrac13+\tfrac1{12}t,&
 k_2^\theta&=\tfrac13+\tfrac1{12}t^{2+\theta},\\
 k_0^\theta&=\tfrac13-\tfrac1{12}(t+t^{2+\theta}).
 \end{aligned}
\end{equation}
These cells are at least $1/6$, sum to one, and span $W_{A_\theta}$.
Use the fixed comparator cube $[1/4,3/4]^3$; every report factor is at
least $\kappa=1/24$. The three fixed probe commands have failure factors
\begin{equation}\label{eq:common-near-probes}
 F_j^\theta=\tfrac12+\tfrac18 k_j^\theta,\qquad j=0,1,2.
\end{equation}
The command accepts cell $j$ with probability $3/8$ and the other cells
with probability $1/2$. Their sum is $13/8$, so these factors span
$W_{A_\theta}$. At a checkpoint $n$ use the $3^{5-n}$ ordered product
queries of length $5-n$. Exactly one checkpoint and one query are used
in a run; all five trials are counted.

For the unconditional experiment the first three commands are independent
uniform vectors in the comparator cube. After their reports have been
compressed, reveal the independent uniform two-step query. The controller
knows $\theta$ and $\mu$ as read-only calibration, but has access to past
commands and reports only through its persistent index. The input law and
all query labels are independent of $\theta$.

For a history write $M_s=\mu(P_h t^s)$ and use the five coordinates
\begin{equation}\label{eq:five-physical-coordinates}
 w=(M_1,M_2,M_{3+\theta},M_{4+2\theta},M_{2+\theta}-M_2).
\end{equation}
They determine every length-two future law. They are ordinary posterior
expectations of physical tests and their differences. For analysis put
\begin{equation}\label{eq:five-desingularized}
 z=(w_1,w_2,w_3,w_4,w_5/\theta)\quad(\theta>0),
 \qquad z_5\big|_{\theta=0}=\mu(P_h t^2\log t).
\end{equation}
This fifth coordinate at zero is an analytic continuation, not an extra
observable supplied to the filter. The zero-calibration future law uses
only the first four coordinates.

\begin{theorem}[Sharp two-parameter streaming order]\label{thm:uniform-streaming}
Set
\begin{equation}\label{eq:five-resolution}
 \Psi_\theta(M)=\max\{M^{-1/2},\theta^{2/5}M^{-2/5}\},
 \qquad M\ge1,
\end{equation}
where the second term is zero at $\theta=0$. There are constants
$0<c\le C<\infty$, depending on the fixed prior and protocol but not on
$M$ or $\theta\in[0,1/2]$, with the following properties.

There is a deterministic $M$-state filter, updated after every report,
whose conditional prediction regret at every checkpoint and every
admitted history is at most $C\Psi_\theta(M)$. In the prescribed
three-trial exploration experiment, every $M$-state filter has
unconditional regret at least $c\Psi_\theta(M)$. The same lower bound
therefore holds for the minimax criterion over histories and checkpoints.
The lower assertion permits the independent randomization of
Definition~\ref{def:finite-state}.
\end{theorem}

In particular the two regimes are
\begin{equation}\label{eq:two-regimes}
 \overline{\mathcal R}^{\theta,\mathrm{stream}}_{M,5}\asymp
 \begin{cases}
 M^{-1/2},&1\le M\le\theta^{-4},\\
 \theta^{2/5}M^{-2/5},&M\ge\theta^{-4},
 \end{cases}
\end{equation}
with uniform comparison constants. This is an order-sharp crossover, not
a claim of an exact transition budget or a leading distortion constant.
The family and its base-model perturbation estimate were introduced in
the source-pinned technical note \cite{A1v6note}. The conclusion here is
the matching, uniform two-parameter law and its online implementation.

\subsection{Attainment of the fifth desingularized direction}
\begin{lemma}[A nondegenerate limiting patch]\label{lem:five-patch}
At a common interior three-command tuple the map to $z$ in
\eqref{eq:five-desingularized} has rank five for every $\theta\in[0,1/2]$.
Its all-failure subprobability law dominates a fixed positive multiple of
uniform probability on a translated five-dimensional cube of fixed side
length, uniformly in $\theta$.
\end{lemma}
\begin{proof}
The formal two-fold sum pairs give exponents
\[
 0,\ 1,\ 2,\ 2+\theta,\ 3+\theta,\ 4+2\theta.
\]
There is just one double cluster, at exponent two. Its divided difference
is $(t^{2+\theta}-t^2)/\theta$. The nonconstant orders are
$(0,0,0,0,1)$. The condition in \eqref{eq:full-future-condition} is
$5\le3(3-1)=6$.

For an explicit common attainable tuple take
$c_1=1/64$, $c_2=2/64$, $c_3=3/64$ and failure factors
$(1+c_it^{2+\theta})/2$. They are realized by rejection probabilities
\[
 (1-g_0,1-g_1,1-g_2)
       =(1/2-2c_i,1/2-2c_i,1/2+4c_i),
\]
strictly inside $[1/4,3/4]^3$. At zero their product tangent is the
entire polynomial space of degree at most six. The future confluent
space has basis
\[
 1,\quad t,\quad t^2,\quad t^2\log t,\quad t^3,\quad t^4.
\]
Lemma~\ref{lem:confluent-positive} makes its mixed pairing with that
tangent surjective. Normalization removes one dimension, not two.
For example, for the uniform prior, the minor with rows $t^i$,
$0\le i\le5$, is exactly
\begin{equation}\label{eq:confluent-exact-minor}
 \det\left[
 \frac1{i+1},\frac1{i+2},\frac1{i+3},
 -\frac1{(i+3)^2},\frac1{i+4},\frac1{i+5}
       \right]_{i=0}^{5}
       =\frac1{338751673344000000}>0.
\end{equation}
This rational identity is only a diagnostic of the preceding all-prior
argument.

For every positive $\theta\le1/2$ the six displayed exponents are distinct
and the binomial tangent has seven distinct exponents. Strict mixed-moment
positivity gives surjectivity throughout this interval. All derivatives
in the desingularized coordinates extend continuously to zero, as proved
in Lemma~\ref{lem:observable-scales}. Compactness and the uniform inverse
construction in Lemma~\ref{lem:uniform-patch} therefore apply on the
\emph{whole} interval $[0,1/2]$, not only on an unspecified punctured
neighborhood. The exact change of order in \eqref{eq:five-desingularized}
is a fixed permutation of the general jet coordinates and changes no
rank or minorization conclusion.
\end{proof}

\subsection{Uniform causal quantization}
\begin{proof}[Proof of Theorem~\ref{thm:uniform-streaming}]
At times one and two retain, exactly for the moment, the chronological
normalized factors $f_i/\mu f_i$, with two coefficient coordinates per
factor. At time three change to $w$ in
\eqref{eq:five-physical-coordinates}. At time four retain
$(M_1,M_{2+\theta})$, and at time five use the unique terminal state.
These states are sufficient for the remaining observable future law.
All their reachable sets are compact.

The first two factor sets lie in fixed boxes of dimensions two and four,
uniformly in $\theta$. Indeed report coefficients are uniformly bounded,
$\mu f\ge\kappa$, and the coefficient hyperplane $\mu f=1$ can be
coordinatized by the two nonconstant coefficients, its constant coefficient
then being $1-\sum_{a>0}f_a\mu(t^a)$. The required prior moments are
fixed read-only data for each known calibration. Uniform boundedness of
these affine formulas on $[0,1/2]$ is immediate.

At time three the first four coordinates of $w$ belong to $[0,1]$.
For $0<t\le1$, the fundamental theorem of calculus in the exponent gives
\[
 |t^{2+\theta}-t^2|
   \le\theta t^2|\log t|\le\theta.
\]
This extends to zero. Consequently $-\theta\le w_5\le0$ for every
posterior, with no evidence denominator in this bound. At time four
both retained moments lie in $[0,1]$. Lemma~\ref{lem:rectangle}, with
reachable representatives in nonempty cells, now gives at most $M$
representatives at each time and squared covering errors
\begin{equation}\label{eq:stage-errors}
 \delta_1^2\le C_1M^{-1},\quad
 \delta_2^2\le C_2M^{-1/2},\quad
 \delta_3^2\le C_3\Psi_\theta(M),\quad
 \delta_4^2\le C_4M^{-1}.
\end{equation}
The $C_i$ are independent of $\theta$. At zero simply omit the last,
constant coordinate of the third state. Every bound in
\eqref{eq:stage-errors} is at most $C_0\Psi_\theta(M)$.

It remains to verify that repeated updates do not divide by $\theta$.
The factor-to-factor update appends the current normalized factor. The
factor-to-$w$ update uses the product of the two retained factors and the
new report factor. It is a ratio of polynomial coefficient expressions
and bounded prior integrals. Along line segments between factor tuples
every interpolated normalized factor is at least $\kappa$, so the evidence
denominators and derivatives have positive uniform bounds. Each physical
test in $w$ is uniformly bounded; differentiating in factor coordinates
therefore gives a uniform Lipschitz constant for this changeover.

For the only nontrivial shrinking update, write the new report factor as
$f=a+bt+ct^{2+\theta}$. In terms of $w$ set
\[
 Z_f=a+bw_1+c(w_2+w_5).
\]
The next two moments are exactly
\begin{equation}\label{eq:uniform-update}
 M'_1=\frac{aw_1+bw_2+cw_3}{Z_f},\qquad
 M'_{2+\theta}=\frac{a(w_2+w_5)+bw_3+cw_4}{Z_f}.
\end{equation}
Coefficients are uniformly bounded and $Z_f\ge\kappa$. A line segment
between reachable $w$ vectors represents a posterior mixture, so the
same lower bound holds along that segment. The quotient rule proves a
uniform Lipschitz estimate in the \emph{physical} $w$ coordinates.
There is no divided difference and no factor $1/\theta$ in
\eqref{eq:uniform-update}. The next, terminal update is constant.

At each time the filter stores only a representative index and applies
the exact transition to that representative and the current input, then
requantizes. A reachable representative followed by any admitted report
has a reachable image, so every quantizer is used inside its domain.
If $e_n$ is the state error in the appropriate Euclidean physical
coordinates, the preceding bounds give
\[
 e_0=0,\qquad e_{n+1}\le L_ne_n+C_0^{1/2}\Psi_\theta(M)^{1/2},
\]
with $L_n$ independent of $\theta$. Iterate this recurrence through the
five times. At times one and two each query prediction is a normalized
polynomial expression in retained factors; at times three and four it
is a fixed linear expression in the retained physical moments. These
prediction maps have uniform Lipschitz constants, again by the positive
evidence bounds. Squaring the resulting prediction error proves the
conditional upper bound at every checkpoint. The representatives and
transition functions are supplied existentially from the exact model,
not synthesized from finite-precision exponent data. Only their index is
persistent mutable memory.

For the lower bound at time three apply
Lemma~\ref{lem:five-patch} and the physical-metric argument in the proof
of Theorem~\ref{thm:confluent-law}, now to the nine common queries in
\eqref{eq:common-near-probes}. Their coefficient change from any other
fixed spanning basis is invertible before taking products, so its full
product expansion has constant full column rank. Four physical directions
have scale one and the last has scale $\theta$. Lemma~\ref{lem:rectangle}
gives the maximum of the four-dimensional lower bound $cM^{-1/2}$ and
the five-dimensional lower bound $c\theta^{2/5}M^{-2/5}$. The constants
are uniform on $[0,1/2]$. Every streaming filter is an $M$-message encoder
at this time, and the retained all-failure evidence is at least $(1/4)^3$.
The randomization argument of Theorem~\ref{thm:confluent-law} applies
unchanged. Finally, minimax regret is at least the unconditional regret
in this one prescribed experiment. This completes the proof.
\end{proof}

\begin{corollary}[Uniform bit law and the lower-constant scale]\label{cor:uniform-bits}
For sufficiently small $\varepsilon>0$, the least persistent bits for
unconditional regret at most $\varepsilon$ in the five-trial protocol,
and likewise for its checkpoint-uniform minimax criterion, satisfy
\begin{equation}\label{eq:uniform-bits}
 B_*(\varepsilon,\theta)
   =\log_2\max\{\varepsilon^{-2},\theta\varepsilon^{-5/2}\}+O(1),
\end{equation}
where the $O(1)$ is uniform in $\theta\in[0,1/2]$. If
\[
 c_{\rm opt}(\theta)=\inf_{M\in\N,\,M\ge1}
           M^{2/5}\overline{\mathcal R}_{M,5}^{\theta,\mathrm{stream}},
\]
then $c_{\rm opt}(\theta)\asymp\theta^{2/5}$ for $0<\theta\le1/2$,
with constants independent of $\theta$.
\end{corollary}
\begin{proof}
Solve the two necessary and sufficient inequalities from
Theorem~\ref{thm:uniform-streaming} for $M$, obtaining, up to constant
factors, $M\ge\varepsilon^{-2}$ and $M\ge\theta\varepsilon^{-5/2}$.
Rounding to $M=2^B$ changes $B$ by a bounded amount. The same bounds
apply to the minimax criterion. The lower bound for $c_{\rm opt}$ follows
from the five-dimensional term. For its upper bound let $M$ tend to
infinity in the matching upper estimate, so
$M^{-1/10}$ disappears. No limiting leading quantization constant is
needed.
\end{proof}

The exact profiles remain $(0,2,4,4,2,0)$ at zero and
$(0,2,4,5,2,0)$ at positive $\theta$. The new assertion is that the extra
direction is resolved at distortion scale $\theta^2$ and state-budget
scale $\theta^{-4}$, with matching estimates on both sides. It sharpens
the base-model bound $C(M^{-1/2}+\theta^2)$ without contradicting it.
